\chapter{A collision surface and its quantum intertwiner}
\label{chap:gi-collision-intertwiner}

Two points can approach one another while retaining a direction of approach.
That direction forms a projective line.
The same projective line appears when two coordinates and their conjugate momenta are reduced by a scalar constraint.
We will identify the surfaces, compute the classes of their distinguished curves, and obtain an operator which exchanges two quantum representations.

The operator will mix two states:
\[
 \frac1{u-\hbar}
 \begin{pmatrix}u&-\hbar\\-\hbar&u\end{pmatrix}.
\]
Here $u$ is the difference of two spectral parameters and $\hbar$ is a deformation parameter.
We will derive both parameters from specified actions and fix their signs.
The off-diagonal entries will come from a curve class, not from a dimension count.

There are three questions to separate.
First, which geometric classes define this matrix?
Second, which algebra acts on the quantum states, and which tensor product does it prescribe?
Third, what remains when the matrix ceases to be invertible?
The last question will distinguish a symmetric quotient from an antisymmetric residue.

\section{The surface that retains a collision direction}

Let $\mathbb P^1$ have affine coordinates $x,x'$ with $x'=x^{-1}$ on the overlap.
Its cotangent bundle $X=T^*\mathbb P^1$ has coordinates $(x,v)$ and $(x',v')$ satisfying
\begin{equation}\label{eq:gi-cotangent}
 x'=x^{-1},\qquad v'=-x^2v.
\end{equation}
Indeed, equality of the covectors $v\,dx=v'\,dx'$ gives this formula.
It preserves $dx\wedge dv$.
The zero section $Z$ consists of the zero covectors.
Write $F_0,F_\infty$ for the cotangent fibres above $0,\infty$.

The same surface describes centered subschemes of length two in the affine plane.
A length-two subscheme is specified by an ideal $I\subset\mathbb C[X_1,X_2]$ whose quotient has dimension two.
For families over a commutative $\mathbb C$-algebra $R$, require this quotient to be locally free of rank two over $R$.
Multiplication by $X_1,X_2$ gives two commuting endomorphisms of that rank-two module.
Their half-traces $c_1,c_2$ are its center coordinates.
Subtract them and denote the resulting trace-zero coordinate classes by $\xi_1,\xi_2$.

\begin{proposition}\label{prop:gi-hilbert-chart}
The space of length-two subschemes of $\mathbb A^2$ is
$\mathbb A^2_{c_1,c_2}\times T^*\mathbb P^1$.
The zero section in its second factor parametrizes the nonreduced centered subschemes.
On the locus of two distinct points, the sum of their symplectic forms extends as
\[
 2\,dc_1\wedge dc_2+dx\wedge dv.
\]
\end{proposition}
\begin{proof}
The trace map splits the quotient module as $R1\oplus L$, where $L$ is locally free of rank one.
The factor $2$ is invertible in $R$.
At every residue field, the trace-zero classes $\xi_1,\xi_2$ generate $L$.
Otherwise both coordinates would be scalar and the quotient algebra would have dimension one.
Thus the open sets where $(1,\xi_1)$ or $(1,\xi_2)$ is a basis cover every family.
This also holds over nonreduced bases: a map onto a finite module is surjective near a point whenever it is surjective over its residue field.
To see this last statement, write the map in local generators and use an invertible minor.

On the first chart, $\xi_2=a\xi_1$.
The trace-zero multiplication matrix of $\xi_1$ satisfies $\xi_1^2=b1$ by its degree-two characteristic equation.
Consequently the ideal is
\[
 (\xi_2-a\xi_1,\ \xi_1^2-b).
\]
Conversely, this quotient is free on $(1,\xi_1)$, and both coordinate multiplications have trace zero.
The second chart has ideal $(\xi_1-c\xi_2,\xi_2^2-d)$.
On the intersection, $c=a^{-1}$ and $d=a^2b$.
Set $x=a$, $v=-b$, $x'=c$, and $v'=d$.
These are exactly \eqref{eq:gi-cotangent}.
The formulas classify families and their changes of basis, so they identify the parameter spaces as varieties, not only their points.
Restoring $c_1,c_2$ gives the product asserted in the proposition.

The centered algebra is nonreduced precisely when $b=0$, or $d=0$ on the other chart.
That is the zero section.
For two distinct centered points, write them as $(r,ar)$ and $(-r,-ar)$, so $b=r^2$.
Their symplectic form is
\[
 2\,dr\wedge d(ar)=db\wedge da=da\wedge d(-b)=dx\wedge dv.
\]
The center contributes $2\,dc_1\wedge dc_2$.
The displayed form is regular and nondegenerate on both charts, including $b=0$.
Thus it supplies the required extension.
\end{proof}

The normal line bundle of $Z$ consists of tangent directions in the surface modulo directions tangent to $Z$.
Here it is the cotangent line bundle of $\mathbb P^1$.
The degree of a line bundle counts zeros minus poles of a nonzero meromorphic section, with multiplicities.
The self-intersection of a compact curve in a surface is the degree of its normal line bundle.
A constant normal section on the first chart has coefficient $-x^2=-(x')^{-2}$ on the second chart.
It therefore has one pole of order two and no zeros.
The normal degree, and hence the self-intersection of $Z$, is $-2$.
We will recover this number from the same equivariant classes that determine the exchange operator.

\section{Equivariant classes from two charts}

Let $T=(\mathbb C^\times)^2$ act on the first chart by
\[
 (s,t)\cdot(x,v)=(s x,s^{-1}t^{-1}v).
\]
Write $u$ for the character of $s$ and $\eta$ for the character of $t$.
The symplectic form has character $-\eta$.
The tangent characters at the two fixed points are
\begin{equation}\label{eq:gi-weights}
 (u,-u-\eta)\quad\hbox{at }0,
 \qquad (-u,u-\eta)\quad\hbox{at }\infty.
\end{equation}
A character such as $u-\eta$ means the one-dimensional representation on which $(s,t)$ acts by $st^{-1}$.
No action on a compact surface containing this local chart is assumed.

We use equivariant cohomology with complex coefficients.
Here is an explicit model and the calculation needed below.
Replace each $\mathbb C^\times$ by its circle subgroup and use the free circle action on $S^\infty$.
For a $T$-space $Y$, its Borel space is
\[
 Y_T=(S^\infty\times S^\infty\times Y)/(S^1\times S^1).
\]
By definition, $H_T^*(Y)$ is the ordinary singular cohomology of this space.
Singular cochains assign complex values to continuous simplices, with differential dual to the alternating sum of faces.
Their cup product evaluates one cochain on the front face and the other on the back face of a simplex.

For a point, the Borel space is $\mathbb P^\infty\times\mathbb P^\infty$.
The space $\mathbb P^N$ has one cell in each even dimension $0,2,\ldots,2N$.
The cellular differentials are zero because there are no odd cells.
The degree-two hyperplane class has its $j$th power evaluating to $1$ on $\mathbb P^j$.
Indeed, $j$ transverse complex hyperplanes meet in one point with positive complex orientation.
These classes give the whole cohomology in the indicated degrees.
Taking $N$ larger than a fixed degree and using the product cells gives
\[
 H_T^*(\mathrm{pt})=S:=\mathbb C[u,\eta],\qquad |u|=|\eta|=2.
\]
The class of a character $s^at^b$ is $au+b\eta$.
This follows from tensoring the corresponding line bundles: their transition winding numbers add.

\begin{lemma}\label{lem:gi-equivariant-ring}
Restriction to the fixed points identifies
\[
 H_T^*(X)=\{(f_0,f_\infty)\in S^2:
                      f_0-f_\infty\in uS\}.
\]
All odd cohomology vanishes.
As an $S$-module, this ring has basis $1=(1,1)$ and $F_0=(u,0)$.
\end{lemma}
\begin{proof}
Multiplying cotangent fibres by a real number between zero and one is a $T$-equivariant deformation to the zero section.
It remains to calculate $H_T^*(\mathbb P^1)$.
The two affine coordinate charts contract equivariantly to their fixed points.
Their intersection is $\mathbb C^\times$, where the first torus factor acts freely and transitively.
Its Borel space is therefore homotopy equivalent to the classifying space of the second factor.
Its cohomology is $\mathbb C[\eta]=S/(u)$, and both restriction maps are the quotient by $u$.

The cohomology sequence for this open cover gives
\[
 0\longrightarrow H_T^{2d}(\mathbb P^1)
 \longrightarrow S^{2d}\oplus S^{2d}
 \xrightarrow{(f,g)\mapsto f-g} (S/(u))^{2d}
 \longrightarrow H_T^{2d+1}(\mathbb P^1)\longrightarrow0.
\]
One can obtain this sequence directly from singular cochains subordinate to the two open sets.
Repeated barycentric subdivision makes each simplex subordinate.
The prism between a simplex and its subdivision gives a chain homotopy to the identity.
Thus subordinate chains calculate the same cohomology, and their short exact cover sequence gives the displayed sequence.
The middle map is surjective in every degree.
Its kernel is the asserted congruence condition, and its cokernel is zero.
Restriction preserves cup products, so this is also the ring description.

Every such pair has the unique expression
\[
 (f_0,f_\infty)=f_\infty(1,1)
          +\frac{f_0-f_\infty}{u}(u,0).
\]
The quotient in this expression is a polynomial by the congruence condition.
This proves the basis statement without inverting $u$.
\end{proof}

A closed smooth complex curve in $X$ has a degree-two class, defined by the Thom class of its normal line.
The Thom class integrates to $1$ on a normal complex disk relative to its boundary.
Restricting it to the curve gives the first Chern class of the normal line.
At a fixed point, that restriction is the normal character.
At a fixed point outside the curve, it is zero.
Thus \eqref{eq:gi-weights} gives
\begin{equation}\label{eq:gi-curves}
 [Z]=(-u-\eta,u-\eta),\qquad
 [F_0]=(u,0),\qquad [F_\infty]=(0,-u).
\end{equation}
The first formula uses the cotangent normal direction.
The other two use the tangent direction of the base.
All three pairs satisfy the congruence in Lemma~\ref{lem:gi-equivariant-ring}.

Integration over the compact zero section sends $1$ to zero for degree reasons and the class of its point $0$ to $1$.
The basis formula therefore gives
\[
 \int_{\mathbb P^1}(f_0,f_\infty)=\frac{f_0-f_\infty}{u}.
\]
Applied to the normal class of $Z$, this yields $-2$.
Equivalently, cancellation of the normal Euler factors gives
\begin{equation}\label{eq:gi-contact}
 \frac{(-u-\eta)^2}{u(-u-\eta)}
 +\frac{(u-\eta)^2}{(-u)(u-\eta)}=-2.
\end{equation}
The equality is first written where the denominators are invertible.
Its value is the integral class $-2$, independent of those inversions.

For the original plane coordinates of the length-two description, let their torus characters be $\alpha,\beta$.
The chart coordinates have characters
\[
 u=\beta-\alpha,\qquad \eta=-(\alpha+\beta).
\]
Indeed, $a=\xi_2/\xi_1$ has character $\beta-\alpha$, and $v=-\xi_1^2$ has character $2\alpha=-u-\eta$.
This identifies the equivariant collision calculation with the cotangent calculation, including the sign of its symplectic character.

\section{Stable classes and their normalization}

Let $A\subset T$ be the first torus factor.
For its positive direction, consider the limit $s\to0$.
A point is attracted to a fixed point if this limit exists and equals that fixed point.
At $0$, the attracting set is the zero section minus $\infty$.
At $\infty$, it is the entire fibre $F_\infty$.
Allowing successive limits gives the full attracting union $Z\cup F_\infty$ for $0$, and $F_\infty$ for $\infty$.
For the opposite direction, interchange $0$ and $\infty$.

A polarization fixes a sign for the product of one character from each opposite pair of $A$-characters.
We choose the cotangent direction.
Its values at the fixed points are $-u$ at $0$ and $u$ at $\infty$.
For either direction, the diagonal normalization is the Euler class of the repelling line, with sign chosen to agree with this polarization after setting $\eta=0$.

In this surface, a stable class for a fixed point means a degree-two equivariant class satisfying three conditions.
It is supported on that point's full attracting union.
Its diagonal restriction has the stated Euler normalization.
Its restriction at a strictly lower fixed point has degree less than one in $u$, with $\eta$ held fixed.
Support means that the class restricts to zero on the complement of the indicated union.

\begin{proposition}\label{prop:gi-stable-classes}
The two directions give the uniquely determined classes
\begin{align*}
 \operatorname{Stab}_+(0)&=[Z]+[F_\infty],&
 \operatorname{Stab}_+(\infty)&=-[F_\infty],\\
 \operatorname{Stab}_-(0)&=-[F_0],&
 \operatorname{Stab}_-(\infty)&=[Z]+[F_0].
\end{align*}
\end{proposition}
\begin{proof}
Each displayed sum has the required support because its summands are classes of curves in that union.
Their restriction matrices, with columns indexed by $0,\infty$, are
\begin{equation}\label{eq:gi-stable-matrices}
 S_+=\begin{pmatrix}-u-\eta&0\\-\eta&u\end{pmatrix},\qquad
 S_-=\begin{pmatrix}-u&-\eta\\0&u-\eta\end{pmatrix}.
\end{equation}
For the positive direction, the repelling line at $0$ has character $-u-\eta$.
At $\infty$ it has character $-u$, whose sign must be reversed to match the cotangent polarization $u$.
These give the two diagonal entries of $S_+$.
For the negative direction, the same rule gives $-u$ and $u-\eta$.
The only nonzero restriction below a diagonal is $-\eta$, of degree zero in $u$.
This proves existence with the stated normalization.

For uniqueness, a class supported on a single fibre has zero restriction at the other fixed point.
The diagonal condition then determines it by the injective restriction map of Lemma~\ref{lem:gi-equivariant-ring}.
For the other class, its lower restriction has degree zero in $u$.
The congruence with the prescribed diagonal restriction determines it uniquely: it must be $-\eta$.
Again the restriction map is injective.
\end{proof}

The stable maps send the two standard generators of $S^2$ to these two classes.
They are maps of $S$-modules, not multiplicative maps for the ordinary cup product.
In the integral basis $(1,[F_0])$ of $H_T^*(X)$, their matrices are
\[
 A_+=\begin{pmatrix}-\eta&u\\-1&-1\end{pmatrix},\qquad
 A_-=\begin{pmatrix}0&u-\eta\\-1&-1\end{pmatrix}.
\]
Their determinants are $u+\eta$ and $u-\eta$.
After inverting $u-\eta$, the change from the positive to the negative stable coordinates is
\begin{equation}\label{eq:gi-geometric-r}
 R_{\mathrm{geom}}(u)=A_-^{-1}A_+
       =\frac1{u-\eta}\begin{pmatrix}u&-\eta\\-\eta&u\end{pmatrix}.
\end{equation}
Thus the formula needs no inverse of $u$.
The fixed-point matrices give the same expression after their apparent extra denominator cancels.
In particular, when $u=0$ and $\eta\ne0$, both stable maps remain invertible in the integral cohomology basis, although fixed-point restriction after specialization no longer separates that basis.

Changing the polarization multiplies each stable column by its chosen sign.
Using the same change in both directions conjugates $R_{\mathrm{geom}}$ by that diagonal sign matrix.
No slope or K\"ahler parameter occurs in these ordinary equivariant cohomology conditions.

\section{The normal critical pair}

The collision surface also arises from a potential with two normal variables.
To retain that comparison, consider coordinates $(x,v,\mu)$ with overlap
\[
 x'=x^{-1},\qquad v'=-x^2v+x\mu,\qquad \mu'=\mu.
\]
Adjoin a global coordinate $\lambda$ and set $W=\lambda\mu$.
The differential of $\mu$ is nonzero on every chart.
Hence the critical equations are $\lambda=\mu=0$, with critical locus $X$.
Let $\mu$ have character $-\eta$ and $\lambda$ character $\eta$.
Then $W$ is invariant under the torus.

We use normal vanishing cohomology to compare cohomological curve operators.
At a point of $X$, take a small normal bidisc $|\mu|,|\lambda|<\delta$.
For $0<|\epsilon|<\delta^2$, its nearby level $W=\epsilon$ is
\[
 \frac{|\epsilon|}{\delta}<|\mu|<\delta,\qquad
 \lambda=\epsilon/\mu.
\]
This is an annulus.
Its reduced cohomology is one-dimensional in degree one and zero in all other degrees.
For the degree-one assertion, contract the annulus radially to a circle and use one zero-cell and one one-cell, with zero boundary.
Choose the class
\[
 \delta_W=\left[\frac1{2\pi i}\frac{d\mu}{\mu}\right],
\]
which integrates to $1$ around the positively oriented circle.

Normalize the vanishing-cohomology complex by shifting this degree-one class to degree zero.
In the convention where the unshifted vanishing-cycle complex has the reduced nearby cohomology as its stalk cohomology, this is the shift $[1]$.
Only ordinary complex cohomology is considered here.
No comparison of Hodge weights is asserted.

\begin{proposition}\label{prop:gi-critical-comparison}
The normalized normal vanishing-cohomology complex is the constant rank-one system on $X$, with generator $\delta_W$.
The map $\delta_W\mapsto1$ identifies its equivariant cohomology with $H_T^*(X)$.
With this normalization at both fixed points, it intertwines the four stable curve classes and their change-of-basis operator.
\end{proposition}
\begin{proof}
Each coordinate chart is a product of the tangent coordinates with the normal pair $(\mu,\lambda)$.
The annulus calculation gives a rank-one local system in normalized degree zero.
On overlaps, $\mu$ and $\lambda$ are unchanged.
Therefore $d\mu/\mu$ represents the same oriented normal class in both charts.
The term $x\mu$ in the tangent transition does not alter it.
More explicitly, replacing this term by $t x\mu$, $0\leq t\leq1$, is an invertible change on the overlap preserving $W$ and fixing $X$.
It connects the transition to the product transition without changing the normal circle.
The local systems consequently glue to the stated constant system.

Multiplication of $\mu$ by a nonzero complex scalar preserves its oriented circle class.
Thus the torus acts trivially on this normal cohomology local system.
The same trivialization applies at the two fixed points.
There is no other stalk cohomology.
A complex with cohomology only in degree zero is identified with that cohomology by the two truncation maps: retain the kernel in degree zero, then quotient by its incoming boundaries.
These maps are equivariant, and give the required quasi-isomorphism of the normalized complex with the constant system.
Cochains on the Borel space with values in this constant local system are ordinary cochains multiplied by $\delta_W$.
This gives the asserted equivariant comparison.
The form $d\mu/\mu$ specifies the ordinary stalk generator.
It is not being used as an equivariantly closed differential form on the whole nearby fibre.

For each curve $C=Z,F_0,F_\infty$, take its Thom class with coefficient $\delta_W$ in that system.
Its support and Euler restrictions are those of $[C]$, multiplied by $\delta_W$.
The four sums in Proposition~\ref{prop:gi-stable-classes} therefore satisfy the same support, normalization, and degree conditions with these coefficients.
Uniqueness follows by the same restriction calculation.
Applying $\delta_W\mapsto1$ gives the four ordinary stable classes, and hence also intertwines their coordinate change.
\end{proof}

This is a comparison of specified cohomology groups and curve operators.
The full derived critical locus also retains tangent directions in nonzero cochain degrees.
The normal calculation does not replace its algebra of functions or its category of matrix factorizations by ordinary cohomology.
The orientation would also change if the two normal coordinates were exchanged: on $\lambda\mu=\epsilon$, $d\lambda/\lambda=-d\mu/\mu$.

\section{The algebra acting on the quantum states}

Let $V$ be the two-dimensional space of polynomials of degree at most one in a variable $q$.
Its basis is $(1,q)$.
The variable $q$ is a state coordinate and is distinct from the equivariant character $u$.
Define
\[
 E=\partial_q,\qquad F=q-q^2\partial_q,\qquad H=1-2q\partial_q.
\]
These operators preserve $V$.
They satisfy $[H,E]=2E$, $[H,F]=-2F$, and $[E,F]=H$.
Their quadratic scalar is $H^2+2H+4FE=3$.
This is the period-one differential-operator module on the projective line: a polynomial of degree at most one is precisely a section of $\mathcal O(1)$ in the affine coordinate.
On the other chart it is represented by $q' f(1/q')$, which is polynomial exactly under this degree bound.

Write $E_{ij}$, $1\leq i,j\leq2$, for the matrix units in this basis.
Thus $E_{12}=E$, $E_{21}=F$, $E_{11}=(1+H)/2$, and $E_{22}=(1-H)/2$.
They satisfy
\[
 [E_{ij},E_{kl}]=\delta_{jk}E_{il}-\delta_{il}E_{kj}.
\]
The Kronecker symbol $\delta_{ij}$ is $1$ if $i=j$ and $0$ otherwise.

To define a tensor product of evaluation representations, we need an algebra with a coproduct.
Work over $\mathbb C[\hbar]$ and introduce formal series
\[
 T_{ij}(w)=\delta_{ij}+\sum_{n\geq1}t_{ij}^{(n)}w^{-n}.
\]
The variable $w$ is an auxiliary spectral variable.
Define the RTT Yangian $\mathsf Y_\hbar$ to be the associative algebra generated by the coefficients, with relations
\begin{equation}\label{eq:gi-rtt-components}
 (w-s)[T_{ij}(w),T_{kl}(s)]
 =\hbar\bigl(T_{kj}(w)T_{il}(s)-T_{kj}(s)T_{il}(w)\bigr).
\end{equation}
The equality means equality of every Laurent coefficient, so it is a collection of algebraic relations.
All tensor products of coefficient algebras below are algebraic.
The formal series organize their coefficients and require no analytic convergence.

\begin{lemma}\label{lem:gi-evaluation-coproduct}
For a central parameter $z$, the rule
\[
 T_{ij}(w)\longmapsto\delta_{ij}+\frac{\hbar E_{ij}}{w-z}
\]
defines a representation on $V$, with $(w-z)^{-1}=\sum_{n\geq0}z^nw^{-n-1}$.
The coproduct is
\[
 \Delta T_{ij}(w)=\sum_a T_{ia}(w)\otimes T_{aj}(w).
\]
It is an algebra homomorphism and is coassociative.
\end{lemma}
\begin{proof}
In the evaluation formula, the left side of \eqref{eq:gi-rtt-components} is
\[
 \frac{\hbar^2(w-s)}{(w-z)(s-z)}
       (\delta_{jk}E_{il}-\delta_{il}E_{kj}).
\]
On the right, the terms quadratic in the matrix units cancel.
The linear terms give the same expression by subtracting $(w-z)^{-1}$ and $(s-z)^{-1}$.
This proves the representation on every coefficient.

For the coproduct, let $P$ exchange the two factors of an auxiliary copy of $\mathbb C^2\otimes\mathbb C^2$.
Write $T_1,T_2$ for $T$ acting on its first or second auxiliary factor.
Multiplying their matrix entries shows that \eqref{eq:gi-rtt-components} is equivalent to
\[
 \bigl(1-\hbar P/(w-s)\bigr)T_1(w)T_2(s)
 =T_2(s)T_1(w)\bigl(1-\hbar P/(w-s)\bigr).
\]
Here the denominators may be cleared before taking coefficients.
For two commuting coefficient algebras, substitute $T(w)=T^{[1]}(w)T^{[2]}(w)$.
Move $T^{[2]}_1$ past $T^{[1]}_2$, apply the displayed relation in the first coefficient algebra, then in the second.
Move $T^{[1]}_1$ past $T^{[2]}_2$ at the end.
The result is the same relation for the matrix product.
Its entries are exactly the proposed coproduct.
Associativity of three matrix factors proves coassociativity.
The counit sends $T(w)$ to the identity matrix.
The formal inverse $T(w)^{-1}$ supplies the antipode, with its coefficients determined recursively by the equation $T(w)T(w)^{-1}=1$.
Inverting the displayed matrix relation verifies the antipode relation in the opposite algebra.
The two matrix inverse identities give the two counit identities.
\end{proof}

Denote this evaluation module by $V(z)$.
The coproduct prescribes the action on $V(z_1)\otimes V(z_2)$.
It retains more information than the diagonal action of the three operators $E,F,H$.
The terms quadratic in matrix units will determine the exchange operator.

\section{The normalized intertwiner}

On $V\otimes V$, let $P(a\otimes b)=b\otimes a$.
For matrix units write $E_{ij}^{(1)}=E_{ij}\otimes1$ and $E_{ij}^{(2)}=1\otimes E_{ij}$.
Multiplying the two evaluation matrices and clearing their denominators gives
\begin{align}
 D_{ij}(w)={}&\delta_{ij}(w-z_1)(w-z_2)
 +\hbar\bigl((w-z_2)E_{ij}^{(1)}+(w-z_1)E_{ij}^{(2)}\bigr)
 +\hbar^2 C_{ij},\label{eq:gi-coproduct-numerator}\\
 C_{ij}={}&\sum_a E_{ia}\otimes E_{aj}=E_{ij}^{(1)}P.\notag
\end{align}
The opposite coproduct has $C_{ij}^{\mathrm{op}}=E_{ij}^{(2)}P$ instead.
The identity for $C_{ij}$ follows by applying both sides to each basis vector.

\begin{theorem}\label{thm:gi-intertwiner}
Set $u=z_2-z_1$.
The operator
\begin{equation}\label{eq:gi-quantum-r}
 R(u)=\frac{uI-\hbar P}{u-\hbar}
\end{equation}
intertwines the coproduct and opposite coproduct representations.
Over $\mathbb C(\hbar,z_1,z_2)$, it is the unique such intertwiner acting as the identity on $1\otimes1$.
\end{theorem}
\begin{proof}
The coefficient of $w$ in \eqref{eq:gi-coproduct-numerator} implies that an intertwiner commutes with every $E_{ij}^{(1)}+E_{ij}^{(2)}$.
Over the stated field, $\hbar$ is invertible.
The diagonal representation splits as the symmetric space of dimension three and the alternating line.
This can be checked directly: repeated application of $F^{(1)}+F^{(2)}$ to $1\otimes1$ gives the three symmetric basis vectors.
The vector $1\otimes q-q\otimes1$ spans a trivial alternating representation.
The diagonal quadratic operator has values $8$ and $0$ on these spaces.
An intertwiner must therefore preserve them.
Commutation with the raising and lowering operators makes it scalar on the symmetric space, and it is scalar on the alternating line.
Hence it has the form $aI+bP$.

For $I+cP$, the difference from the required intertwining identity is
\[
 (I+cP)D_{ij}-D_{ij}^{\mathrm{op}}(I+cP)
 =\bigl(\hbar^2-\hbar c(z_1-z_2)\bigr)
                     (E_{ij}^{(1)}-E_{ij}^{(2)})P.
\]
To derive it, use $PE_{ij}^{(1)}=E_{ij}^{(2)}P$ and $P^2=I$ in each term of \eqref{eq:gi-coproduct-numerator}.
The quadratic terms proportional to $c$ cancel.
Thus $c=\hbar/(z_1-z_2)=-\hbar/u$.
Dividing by the symmetric eigenvalue $1-\hbar/u$ gives \eqref{eq:gi-quantum-r}.
Equivalently, solving for $a,b$ with $a+b=1$ gives $a=u/(u-\hbar)$ and $b=-\hbar/(u-\hbar)$.
This proves uniqueness over the stated field and the identity itself after clearing denominators.
Consequently the formula specializes wherever $u-\hbar$ is invertible.
At $\hbar=0$, the evaluation matrices become identities, so generic uniqueness is not a claim about that degenerate representation.
\end{proof}

To compare with geometry, include the three varieties $T^*\operatorname{Gr}(k,2)$ for $k=0,1,2$.
Here $\operatorname{Gr}(k,2)$ is the space of $k$-dimensional subspaces of $\mathbb C^2$.
For $k=0,2$ it is a point, and for $k=1$ it is $\mathbb P^1$.
Its fixed coordinate subspaces are indexed by subsets of $\{1,2\}$ of size $k$.
Identify a tensor basis vector with that subset by recording the positions occupied by $q$ rather than $1$.
Thus $q\otimes1$ corresponds to the line spanned by the first framing vector, and $1\otimes q$ to the second.

The stable maps on the two point components are identities.
On the middle component, Proposition~\ref{prop:gi-stable-classes} gives \eqref{eq:gi-geometric-r} in exactly this label order.
Apply the coefficient map
\[
 \mathbb C[u,\eta]\longrightarrow\mathbb C[\hbar,z_1,z_2],
 \qquad u\longmapsto z_2-z_1,\qquad\eta\longmapsto\hbar.
\]
Then the direct sum of the geometric operators is \eqref{eq:gi-quantum-r} on all four tensor states.
The normal vanishing-cohomology comparison of Proposition~\ref{prop:gi-critical-comparison} preserves the same operator with the specified normal orientation.
This gives the geometric and quantum comparison on actual representations, with explicit source, target, basis, and coefficient map.

The fixed-period differential-operator quotient itself does not carry the inherited enveloping coproduct.
Its single-site quadratic value is $3$, whereas the diagonal tensor action has values $8$ and $0$.
The tensor representations above use the full RTT algebra.
They do not impose the single-site quadratic equation on the total tensor action.

\section{Three factors and the singular parameters}

The comparison must be consistent when three factors are reordered.
This requires a polynomial identity that can be proved without a dimension restriction.
Let $u=z_2-z_1$ and $\upsilon=z_3-z_2$ be the two adjacent spectral differences.
For any vector space, let $A=P_{12}$, $B=P_{23}$, and $C=P_{13}$ on its third tensor power.
Their permutation relations are
$A^2=B^2=1$ and $ABA=BAB=C$.

\begin{proposition}\label{prop:gi-yang-baxter}
Where the scalar denominators are invertible,
\[
 R_{12}(u)R_{13}(u+\upsilon)R_{23}(\upsilon)
 =R_{23}(\upsilon)R_{13}(u+\upsilon)R_{12}(u),\qquad
 R(u)R(-u)=I.
\]
These identities hold on every third tensor power, not only on the two-dimensional evaluation module.
\end{proposition}
\begin{proof}
Clear the common denominator $(u-\hbar)(u+\upsilon-\hbar)(\upsilon-\hbar)$.
The first identity becomes equality of
\[
 (uI-\hbar A)((u+\upsilon)I-\hbar C)(\upsilon I-\hbar B)
\]
with its reverse order.
The constant and linear coefficients in $\hbar$ agree term by term.
The quadratic coefficient on the first side is
$uCB+(u+\upsilon)AB+\upsilon AC$.
Since $CB=AC=BA$, it is $(u+\upsilon)(AB+BA)$.
On the other side, $CA=BC=AB$ gives the same result.
The cubic terms agree because $ACB=BCA=C$.
This proves the polynomial identity in the permutation algebra, hence on every tensor space.
For the second identity, $P^2=I$ gives numerator $u^2-\hbar^2$ and the same scalar denominator after cancelling signs.
\end{proof}

For an ordered list of evaluation parameters, the operator $PR(z_{i+1}-z_i)$ maps the corresponding tensor representation to the representation with those parameters exchanged.
The permutation changes the opposite coproduct back into the coproduct in the exchanged order.
Operators acting on disjoint pairs commute.
The preceding identities give the two braid relations and the inverse relation, with parameters exchanged at each step.
Any two adjacent-transposition expressions for the same permutation differ by these relations: move the last label to its prescribed position and induct on the remaining labels.
Thus the resulting representation map is independent of the expression wherever all required pairwise factors $z_i-z_j\pm\hbar$ are invertible.
This proves the algebraic coherence of the tensor comparison.
A geometric factorization statement for larger Grassmannians requires its own fixed-locus and stable-class construction.

The spectral singularities have a direct representation meaning.
Let
\[
 \Pi_+=\tfrac12(I+P),\qquad \Pi_-=\tfrac12(I-P)
\]
be the symmetric and alternating projectors.
Then
\begin{equation}\label{eq:gi-channels}
 R(u)=\Pi_++\frac{u+\hbar}{u-\hbar}\Pi_-.
\end{equation}
For $\hbar\ne0$, its value at $u=-\hbar$ is $\Pi_+$, while its residue at $u=\hbar$ is $2\hbar\Pi_-$.
The latter means the value of $(u-\hbar)R(u)$ at that point.
Both follow directly from \eqref{eq:gi-channels}.
The polynomial intertwining identity shows that the corresponding image and kernel are invariant under the appropriate source or target representation.
These are actual quotient and subrepresentation operations at the singular parameters.

At $u=0$, $\hbar\ne0$, the operator equals $P$.
At $\hbar=0$, $u\ne0$, it equals $I$.
There is no common value at $(u,\hbar)=(0,0)$ independent of the approach: along $u=c\hbar$, the alternating eigenvalue is $(c+1)/(c-1)$.
Thus a simultaneous zero-parameter limit requires additional data.

Give $(1,q)$ its standard Hermitian inner product.
Then $P$ and both projectors are self-adjoint.
For real $\hbar$ and purely imaginary $u$, with $u-\hbar\ne0$, the alternating eigenvalue in \eqref{eq:gi-channels} has absolute value one.
Hence $R(u)$ is unitary in that regime.
For general real $u,\hbar$ it need not be unitary, even though the rational inverse identity still holds.
The symmetric and alternating channels thus give a precise two-state quantum interpretation without identifying the formal spectral variable with a spacetime coordinate.

\section{Surface geometry and threefold geometry}

The cotangent surface used above is a holomorphic symplectic variety.
A smooth complex threefold cannot itself carry a nondegenerate holomorphic alternating two-form.
At any point such a form would be an invertible $3\times3$ skew matrix $M$.
But $\det M=\det(-M^{\mathsf T})=-\det M$, so $\det M=0$ over $\mathbb C$.
In particular, the threefold $\operatorname{Tot}(\mathcal O(-1)\oplus\mathcal O(-1)\to\mathbb P^1)$ is not this cotangent surface.
Using a surface chart in a threefold problem requires the chart and its normal comparison to be specified.
The potential $W=\lambda\mu$ provides such a normal comparison at the stated cohomological level.

The same distinction applies to changes of coefficients or constructions.
Our stable classes are ordinary equivariant cohomology classes, determined by a chamber and polarization.
They produce an ordinary rational Yang--Baxter operator.
A parameter-dependent change of basis is not, by itself, a dynamical Yang--Baxter equation.
Such an equation includes specified shifts of an additional parameter during composition.
No such shifts occur in the identities proved here.

The concrete comparison has retained the collision direction, its exceptional curve, the normal vanishing class, and the full tensor intertwiner.
Its nonzero mixing coefficient and its two singular channels come from these specified structures.
Extending it to larger geometric or categorical objects requires maps preserving those structures, rather than an equality of their dimensions or scalar invariants.
